\documentclass[12pt]{article}

\usepackage[a4paper,margin=2.5cm]{geometry}
\usepackage[T1]{fontenc}
\usepackage[utf8]{inputenc}
\usepackage{lmodern}
\usepackage{microtype}
\usepackage{setspace}
\usepackage{graphicx}
\usepackage{float}
\usepackage{booktabs}
\usepackage{array}
\usepackage{enumitem}
\usepackage{xcolor}
\usepackage{listings}
\usepackage[hidelinks]{hyperref}

\definecolor{codegray}{RGB}{245,245,245}
\definecolor{rulegray}{RGB}{190,190,190}

\setstretch{1.12}
\setlength{\parindent}{0pt}
\setlength{\parskip}{0.6em}
\setlist{itemsep=0.25em,topsep=0.35em}
\renewcommand{\arraystretch}{1.25}
\graphicspath{{resources/}}

\lstset{
    basicstyle=\ttfamily\small,
    backgroundcolor=\color{codegray},
    frame=single,
    rulecolor=\color{rulegray},
    breaklines=true,
    showstringspaces=false,
    columns=fullflexible,
    keepspaces=true,
    aboveskip=0.8em,
    belowskip=0.8em
}

\newcommand{\commandname}[1]{\texttt{#1}}
\newcommand{\labimage}[3][0.92\textwidth]{%
    \begin{figure}[H]
        \centering
        \includegraphics[width=#1,height=0.72\textheight,keepaspectratio]{\detokenize{#2}}
        \caption{#3}
    \end{figure}
}

\begin{document}

\pagenumbering{gobble}

\begin{titlepage}
    \centering

    {\large Ministerul Educa�iei �i Cercet�rii al Republicii Moldova\par}
    \vspace{0.45cm}
    {\large Universitatea Tehnic� a Moldovei\par}
    \vspace{0.45cm}
    {\large Facultatea Calculatoare, Informatic� �i Microelectronic�\par}

    \vspace{3cm}

    {\Large\bfseries Laboratory Work 1\par}
    \vspace{0.55cm}
    {\LARGE\bfseries Introduction to Linux and Operating System Resources\par}
    \vspace{0.35cm}
    {\large Operating Systems\par}

    \vfill

    \begin{flushleft}
        \textbf{Elaborated by:}\\
        Student, group FAF-241\\
        Ca� Chiril
    \end{flushleft}

    \vfill

    {Chi�in�u -- 2026\par}
\end{titlepage}

\pagenumbering{roman}
\tableofcontents
\clearpage
\pagenumbering{arabic}

\section{Introduction}

An operating system manages the computer's hardware and provides abstractions that programs can use safely. Its main responsibilities include organizing files, controlling processes, allocating memory, managing persistent storage, and exposing input/output devices. Linux makes these responsibilities observable through ordinary command-line utilities and virtual filesystems such as \commandname{/proc} and \commandname{/dev}.

This laboratory work examines those facilities from inside a Docker container. The container provides an isolated Linux userspace while sharing the host's kernel. Consequently, the commands show the container's process namespace, filesystem view, and assigned resources rather than a complete view of the macOS host.

\subsection{Objective}

The objective is to become familiar with the Linux command line and to observe how the operating system represents and manages files, processes, memory, storage, mounted filesystems, and devices.

\subsection{Tasks}

The practical work was divided into four parts:

\begin{enumerate}
    \item inspect the home directory and file metadata, then change file permissions;
    \item create and inspect a background process and interpret its state;
    \item examine system memory and the resident memory used by one process;
    \item inspect block devices, the root mount, and device files.
\end{enumerate}

\subsection{Environment and Method}

All commands were executed as the \commandname{root} user in a Docker container. Each command was run interactively, its output was captured, and the result was interpreted in relation to the resource being managed. Table~\ref{tab:commands} summarizes the principal tools used.

\begin{table}[H]
    \centering
    \caption{Commands used during the laboratory work}
    \label{tab:commands}
    \begin{tabular}{>{\ttfamily}p{0.25\textwidth} p{0.62\textwidth}}
        \toprule
        \normalfont\textbf{Command} & \textbf{Purpose} \\
        \midrule
        ls, chmod & Inspect directory entries and modify permission bits. \\
        jobs, ps & Display shell jobs and processes visible in the container. \\
        /proc/PID/status & Read kernel-maintained status and memory data for a process. \\
        free & Summarize physical memory and swap usage. \\
        lsblk, findmnt & Inspect block devices and mounted filesystems. \\
        ls /dev & List device nodes exposed to the container. \\
        \bottomrule
    \end{tabular}
\end{table}

\section{Filesystem and File Permissions}

\subsection{Inspecting the Home Directory}

The command \commandname{ls -la \textasciitilde} lists every entry in the home directory, including hidden files whose names begin with a dot. The long format reports the file type and permissions, link count, owner, group, size, modification time, and name. In the container, the home directory is \commandname{/root}; it contains the standard shell configuration files \commandname{.bashrc} and \commandname{.profile}.

\labimage[0.78\textwidth]{resources/ls -la ~ output.png}{Long listing of the root user's home directory}

The owner and group fields are both \commandname{root}, which is expected because the commands were executed by the container's root account. The first character of the mode identifies the object type: \commandname{d} denotes a directory and \commandname{-} denotes a regular file. The remaining nine characters form three permission groups for the owner, group, and others.

\subsection{Changing Permissions}

Before the permission change, \commandname{note.txt} had the mode \commandname{-rw-r--r--}. This means that the owner could read and write the file, while the group and all other users could only read it.

\labimage[0.72\textwidth]{resources/ls -l note.txt before chmod.png}{Permissions of note.txt before applying chmod}

The command \commandname{chmod 600 note.txt} assigns read and write permission to the owner and removes every permission from the group and others. The octal value is interpreted as (6=4+2) for owner read and write, followed by zero permissions for the other two classes.

\labimage[0.72\textwidth]{resources/ls -l note.txt after chmod.png}{Permissions of note.txt after applying chmod 600}

After the change, the mode is \commandname{-rw-------}. The initial dash still marks a regular file; \commandname{rw-} grants the owner read and write access, and the two \commandname{---} groups deny access to the group and others. This demonstrates how Linux enforces discretionary access control through per-file ownership and permission bits.

\subsection{Purpose of Important Directories}

Two important top-level directories are:

\begin{itemize}
    \item \commandname{/etc}, which stores system-wide configuration files;
    \item \commandname{/var}, which stores variable data that changes while the system runs, including logs, caches, spool data, and some container-runtime storage.
\end{itemize}

\section{Processes and Process States}

\subsection{Creating a Background Process}

The command \commandname{sleep 300 \&} starts \commandname{sleep} as a background job. The shell immediately returns a prompt and reports the job number and process identifier. The \commandname{jobs} command shows that the job is running, while \commandname{ps -ef | grep sleep} locates it in the process table.

\labimage[0.88\textwidth]{resources/ps aux.png}{Starting and locating a background sleep process}

A background job is still an ordinary process. The ampersand only instructs the shell not to wait synchronously for its completion. This permits other commands to be entered while \commandname{sleep} remains active.

\subsection{Process Listing and PID 1}

The \commandname{ps aux} output includes the user, PID, CPU and memory percentages, virtual and resident memory, terminal, state, start time, accumulated CPU time, and command. In this container, PID~1 is \commandname{bash}, as shown in Figure~\ref{fig:pid1}. In a full Linux system PID~1 is normally an init system, but a container can launch its main application directly as PID~1 in its isolated PID namespace.

\begin{figure}[H]
    \centering
    \includegraphics[width=0.9\textwidth]{\detokenize{resources/process with pid 1 is bash.png}}
    \caption{Process list showing bash as PID 1}
    \label{fig:pid1}
\end{figure}

The command \commandname{ps aux | wc -l} returned 5. This value counts the header and the short-lived commands that form the pipeline as well as the persistent processes. It therefore describes the lines visible during that particular measurement, not a permanent process count.

\labimage[0.64\textwidth]{resources/number of processes.png}{Counting the lines produced by ps aux}

\subsection{Reading Process Information from proc}

Linux exposes process information through \commandname{/proc}. After saving the process identifier in \commandname{PID}, the command \commandname{cat /proc/\$PID/status | head -20} displays the beginning of the kernel-maintained status record. The record identifies the command as \commandname{sleep}, gives PID~40 and parent PID~1, and reports its credentials and memory information.

\labimage[0.88\textwidth]{resources/cat :proc:$PID:status | head -20.png}{Beginning of the status record for the sleep process}

The \commandname{State} line contains \commandname{S (sleeping)}. This is the expected state because \commandname{sleep} is waiting for a timer rather than actively using the CPU. A sleeping process remains known to the kernel and can become runnable when the awaited event occurs.

\labimage[0.88\textwidth]{resources/the state line S for the slepping process.png}{Status output showing the process in the interruptible sleeping state}

Together, \commandname{ps} and \commandname{/proc} demonstrate two views of the same kernel data: \commandname{ps} formats a convenient process table, whereas \commandname{/proc/PID} exposes detailed per-process information as virtual files.

\section{Memory Management}

\subsection{System Memory and Swap}

The command \commandname{free -h} presents memory values in human-readable units. At the time of observation, the environment reported 7.7~GiB of physical memory, with approximately 4.1~GiB used, 1.3~GiB free, 2.6~GiB in buffers/cache, and 3.6~GiB available. It also reported 1.0~GiB of swap, unused at the time of measurement.

\labimage[0.95\textwidth]{resources/free -h result.png}{Human-readable physical memory and swap statistics}

The \commandname{free} and \commandname{available} columns have different meanings. \commandname{free} counts memory that currently holds no data, while \commandname{available} estimates how much memory could be supplied to new applications without swapping, including reclaimable cache. Linux deliberately uses otherwise idle memory for caching, so a comparatively small \commandname{free} value does not by itself indicate a memory shortage.

Swap is disk-backed space that can hold memory pages when the kernel needs to reduce pressure on physical RAM. It is much slower than RAM and should be regarded as a fallback mechanism, not additional memory with equivalent performance. The observed zero swap usage indicates that no pages had been moved to swap at that moment.

\subsection{Memory Used by the sleep Process}

A longer \commandname{sleep 1000} process was started and its PID was saved using the shell variable \commandname{pid=\$!}. Querying \commandname{VmRSS} in \commandname{/proc/\$pid/status} reported 1284~kB.

\labimage[0.78\textwidth]{resources/sleep 1000 & vm rss.png}{Resident memory reported for the sleep process}

\commandname{VmRSS} is the resident set size: the portion of a process's mapped memory that is currently present in physical RAM. Although \commandname{sleep} performs almost no active computation, it is still a user-space program and requires executable code, runtime libraries, stack pages, and bookkeeping data. A resident size of about 1.3~MB is therefore reasonable and remains very small relative to the available system memory.

\section{Storage, Mounts, and Devices}

\subsection{Block Devices}

The \commandname{lsblk} utility lists block devices and their relationships. The output shows several network block devices and two virtual disks. Device \commandname{vda} has a partition \commandname{vda1}, while entries associated with \commandname{vdb} are mounted at container configuration files such as \commandname{/etc/resolv.conf}, \commandname{/etc/hosts}, and \commandname{/etc/hostname}.

\labimage[0.86\textwidth]{resources/lsblk output.png}{Block devices visible from inside the Docker container}

The device view is virtualized by Docker. It should not be interpreted as a direct inventory of the Mac's physical disks; it is the storage topology exposed to the Linux environment.

\subsection{The Root Filesystem}

The command \commandname{findmnt /} reports the filesystem mounted at the container root. The result identifies the source as \commandname{overlay} and the filesystem type as \commandname{overlay}.

\labimage[0.92\textwidth]{resources/findmnt : output.png}{Mount information for the container root filesystem}

OverlayFS combines multiple directory layers into one filesystem view. Docker uses this mechanism so that read-only image layers can be shared while each container receives a writable upper layer. Paths in the mount options refer to Docker's internal storage, which in Docker Desktop ultimately resides in the Linux virtual machine backed by the macOS host. Therefore, the result does not mean that the Linux root is simply the host's \commandname{/var} directory; rather, the container root is an overlay assembled from runtime-managed layers.

\subsection{Device Nodes}

The \commandname{ls /dev} command lists the device nodes made available to the container. The result includes terminal-related nodes such as \commandname{console}, \commandname{tty}, \commandname{ptmx}, and \commandname{pts}; standard-stream links such as \commandname{stdin}, \commandname{stdout}, and \commandname{stderr}; and special character devices such as \commandname{null}, \commandname{zero}, \commandname{random}, and \commandname{urandom}.

\labimage[0.98\textwidth]{resources/ls :dev i.e all the devices.png}{Device nodes exposed in the container's dev directory}

Device nodes provide a file-oriented interface to kernel drivers and pseudo-devices. For example, writes to \commandname{/dev/null} are discarded, reads from \commandname{/dev/zero} return zero bytes, and \commandname{/dev/urandom} supplies pseudorandom data. The limited list is another effect of container isolation: Docker exposes only the devices required or explicitly permitted for the container.

\section{Discussion}

The observations connect the principal responsibilities of an operating system to concrete Linux interfaces:

\begin{itemize}
    \item \textbf{Files and protection:} directory listings expose metadata, and \commandname{chmod} changes access rights through permission bits.
    \item \textbf{Processes and CPU scheduling:} \commandname{ps}, shell job control, and \commandname{/proc} reveal process identities, parent-child relationships, and execution states.
    \item \textbf{Memory:} \commandname{free} summarizes system-wide allocation, while \commandname{VmRSS} reports resident memory for one process.
    \item \textbf{Storage and filesystems:} \commandname{lsblk} displays block devices, and \commandname{findmnt} identifies the overlay filesystem used for the container root.
    \item \textbf{Input and output:} \commandname{/dev} presents devices and pseudo-devices through filesystem entries.
\end{itemize}

Containerization affects every observation. The container shares the host kernel but operates with isolated namespaces and a restricted resource view. PID~1 is the container's main shell, the root filesystem is an overlay, and only selected devices are present. These differences are useful examples of how the kernel can present distinct operating-system environments to different groups of processes.

\section{Conclusion}

The laboratory work demonstrated how Linux exposes the resources managed by an operating system through a small set of command-line tools and virtual filesystems. File listings and \commandname{chmod} showed ownership and access control; job-control commands, \commandname{ps}, and \commandname{/proc} showed process creation and state; \commandname{free} and \commandname{VmRSS} described memory at system and process level; and \commandname{lsblk}, \commandname{findmnt}, and \commandname{/dev} revealed the container's storage and device interfaces.

The results also showed why the execution environment matters. Because the commands ran in Docker, they describe an isolated Linux container rather than the complete macOS host. Even so, the same core abstractions---files, processes, virtual memory, mounted filesystems, and device nodes---are the mechanisms through which Linux organizes and protects computing resources.

\end{document}
